\documentclass{article}
\usepackage{amsmath}
\begin{document}
\title{A non-contact cave-mouth read trades earlier evidence for a later return}
\maketitle

\paragraph{Source and kernel.}
Pinned SubT Cave01 places tile 13 at \((100,50,0)\) and 42
at \((100,75,0)\). The pinned connection catalogue gives both
straight models the local ports at \((0,\pm12.5,0)\).
Transforming the declared SDF poses gives the unique matching
mouth at \((100,62.5,0)\), with discrepancy below 0.035 mm
between the two transformed port positions. This is a design
anchor, not an aisle envelope: the graph helper matches
nominal port positions and does not certify the collision mesh,
vehicle footprint, line of sight, or a safe sensor read.
The publicly listed Fuel assets are not part of this replay.

On that axis the authored development fixture uses \(x=0\)
for the near site, \(x=25\) for the far site, and a transverse
barrier at \(12.5\) in the closed branch. A rectangular clear
tube is stipulated with lateral half width 1.3 and sufficient
floor and ceiling clearance. The vehicle is constrained to
this axis, and its front extent plus required gap is 0.5, so
the closed center limit is \(S=12\). The open branch has
a positive optical target at \(x=14\) on a side panel
with lateral support \([0.7,1.2]\), outside the swept body.
The fixture admits one read at a designated marked cut \(d\),
not arbitrary early use of raw optical history. At \(d=8\)
an off-axis ray with the stipulated pose and slope bounds
has a clear beam tube before the closed barrier or open target.
The nominal first-return distances separate even with
0.05 m range error and remain inside 6.2 m.
Neither a missing return nor a collision is used as the
observation. Actual frame exposure time and motion safety in
the original asset remain unverified.

\paragraph{Exact curve for the declared drive.}
Start from rest and stay within the corridor. Absolute
acceleration is bounded by \(a\) and speed by \(V\). A permitted
instantaneous read at \(d\), while moving toward the gate,
makes the branch available to the local installer immediately.
The observer's radio/decoder delay to a different actor is a
separate nonnegative budget. Let \(L\) be the far distance
and assume \(d\le S<L\). To leave stopping room in the closed
branch and reach the read with forward speed \(v\), exactly
the following nonnegative speeds are available:
\[
v\le v_*=\min\{V,\sqrt{2ad},\sqrt{2a(S-d)}\}.
\]
The far branch must stop at \(L\); the closed branch must
return to 0 at rest before placing the near relay.
With \(T(z,v)\) the minimum time for a leg of length \(z\)
from rest to speed \(v\) in the same direction, the best
times at the read, far site, and near return are
\[
R(v)=T(d,v),\quad
F(v)=T(d,v)+T(L-d,v),\quad
N(v)=T(d,v)+v/a+T(d+v^2/(2a),0).
\]
These are achieved using a common prefix.
For \(z\ge v^2/(2a)\), put
\(w=\min(V,\sqrt{az+v^2/2})\), then the optimal leg
accelerates to \(w\), coasts as needed, and brakes to \(v\).
For a given duration its speed is bounded pointwise
by the minimum of the ramp from the initial velocity, the cap,
and the backwards ramp from the destination velocity.
The displayed schedule attains this envelope and hence the
distance at its minimum time. Time reversal handles the far
suffix. In the closed branch maximal immediate braking
gives the earliest turnaround and smallest forward overrun;
any later turnaround has weakly larger time and return
distance. The abstract maneuver must stay short of \(S\).

Through the uncapped segment, \(T(z,v)+v/a=2\sqrt{az+v^2/2}/a\);
through the capped segment it is \(z/V+V/a+v^2/(2aV)\).
The branches agree at the transition. On the feasible interval
\(R,F\) are nonincreasing and \(N\) is nondecreasing.
The first two deadlines impose lower speed bounds and
the return deadline an upper one. Their interval intersection
is the exact feasibility test for this maneuver class.
Radio/decoder latency subtracts its justified bound from the
recipient's decision cutoff; stage time subtracts from the
corresponding relay-ready deadline. They do not follow from a
positive range test. If an old receipt is required before
activation, its availability must additionally precede that
stage cut.

\paragraph{Positive and impossible arms.}
At \(a=1,V=3,L=25,d=8,S=12\), speed 2 at the mark gives
read \(13/3\), far readiness \(35/3\), near readiness \(38/3\)
seconds. Thus read by 5, far by 12 and near by 13 is
feasible, and stopping for the same short-range read misses
the read cutoff. Read by 5 with both sites ready by 12 is
impossible even with instant local response and perfect radio.
At \(v=\sqrt2\) the return is 12 and the far
arrival exceeds 12. At any larger \(v\) the return exceeds 12
(strict monotonicity there); at any smaller \(v\) the far
arrival cannot improve. A truly independent earlier
stand-off observation or an already available far relay might
change the task; it was not ruled out or given to one
comparison arm for free. No actual 72 m delivery time,
robot dynamics, full-cave clearance, or independent mission
is inferred. The output is a quantified replacement for the
previous model's abstract safe failed traversal, under explicit
authored geometry and ideal local sensing.
\end{document}
